\documentclass[conference]{IEEEtran}
\usepackage[utf8]{inputenc}
\usepackage{amsmath}
\usepackage{array}

\title{Metodología para Entrenamiento y Test de un Agente de Navegación en ROS 2}
\author{\IEEEauthorblockN{Grupo de Robótica IMT-342}}

\begin{document}
\maketitle

\section{Configuración del Entorno}
El laboratorio se plantea sobre ROS 2 Jazzy, usando un robot móvil diferencial simulado. El entorno recomendado es Gazebo, con visualización en RViz2. El robot utiliza LiDAR 2D en \texttt{/scan}, odometría en \texttt{/odom} y actuación mediante \texttt{/cmd\_vel}.

\section{Definición del Agente}
El agente recibe lecturas LiDAR reducidas, distancia a la meta, ángulo relativo, velocidad lineal, velocidad angular y distancia mínima a obstáculos. Las acciones son \texttt{linear.x} en $[0.0,0.35]$ m/s y \texttt{angular.z} en $[-1.0,1.0]$ rad/s. Se recomienda PPO.

\[
r_t = 2.0(d_{t-1}-d_t)+0.5\cos(\theta_{goal})+0.1v_x-0.05|w_z|-0.01-P_c+R_g
\]

donde $P_c=100$ ante colisión y $R_g=100$ al alcanzar la meta.

\section{Entrenamiento en Simulación}
El entrenamiento se realiza exclusivamente en simulación mediante episodios con metas aleatorias. Se registra estado, acción, recompensa y condición terminal. Se usa semilla fija 42, máximo 1000 episodios, 500 pasos, evaluación cada 25 episodios y parada temprana tras 10 evaluaciones sin mejora.

\section{Evaluación y Pruebas}
La evaluación se realiza con el agente entrenado usando política determinista. Se ejecutan entre 20 y 30 episodios con metas fijas y aleatorias, registrando éxito, colisión, recompensa acumulada, pasos, tiempo y distancia final.

\section{Métricas de Validación}
\begin{table}[h]
\centering
\begin{tabular}{|p{2.2cm}|p{2.6cm}|p{2.2cm}|}
\hline
Métrica & Descripción & Criterio \\
\hline
Recompensa & Suma por episodio & Creciente \\
Éxito & Metas alcanzadas & $\geq 80\%$ \\
Colisión & Episodios con choque & $\leq 10\%$ \\
Tiempo & Duración media & Decreciente \\
Distancia final & Error a la meta & $<0.30$m \\
Pasos & Acciones ejecutadas & Decreciente \\
\hline
\end{tabular}
\end{table}

\section{Conclusión}
La metodología permite entrenar y evaluar un agente de navegación en ROS 2 Jazzy usando LiDAR, odometría y comandos \texttt{/cmd\_vel}. La recompensa integra avance, orientación, seguridad y eficiencia.

\end{document}
